% Generated draft: requires mathematical and source review.
\documentclass[11pt]{article}
\usepackage[margin=25mm]{geometry}
\usepackage{fontspec}
\setmainfont{Latin Modern Roman}
\usepackage{amsmath,amssymb,mathtools,unicode-math}
\setmathfont{Latin Modern Math}
\usepackage{enumitem,needspace,xurl,hyperref}
\setlength{\parindent}{0pt}
\setlength{\parskip}{5pt}
\setlength{\emergencystretch}{3em}
\newcommand{\N}{\mathbb{N}}
\newcommand{\Z}{\mathbb{Z}}
\newcommand{\Q}{\mathbb{Q}}
\newcommand{\R}{\mathbb{R}}
\newcommand{\C}{\mathbb{C}}
\newcommand{\F}{\mathbb{F}}
\newcommand{\A}{\mathbb{A}}
\newcommand{\PP}{\mathbb{P}}
\newcommand{\E}{\mathbb{E}}
\newcommand{\eps}{\varepsilon}
\newcommand{\dd}{\,\mathrm{d}}
\begin{document}
\section*{UnsolvedMath ID: 20002485}
\textbf{Source problem code:} catalogue entry
\subsection{Definitions and mathematical statement}
The canonical extract concatenates two questions. The first, attributed to Franz Auer, conjectures a Hausdorff-dimension bound $n-7$ for the singularity set of a minimizing closed codimension-one current. The second defines a stable norm on $H_q(M;\R)$ by minimum mass and asks for convexity and differentiability properties of its unit ball, comparing it with Mather's $\beta$-function. Both targets are retained.

For the first question, the aid uses ambient dimension $N$ and real codimension-one currents $T$ on a smooth Riemannian $N$-manifold. ``Closed'' means $\partial T=0$; the precise local or homological mass-minimizing class and the definition of singular set must be fixed. The original intended assertion is
\[
 \dim_H\operatorname{sing}(T)\le N-7.
\]
The aid reports a stronger, specifically Auer--Bangert aggregate-singular-set assertion
\[
 \dim_H\operatorname{sing}_{\mathrm{AB}}(T)\le N-8.
\]
Here the aggregate is associated with the measured-leaf decomposition of a minimizing real current, rather than an unspecified notion of singular support. The supplied record does not give its exact definition or all decomposition hypotheses. It attributes the claimed bound to a disjoint-family Hausdorff-content argument of Chodosh--Mantoulidis--Schulze, a $C^2$ variable-metric clustering extension, the Auer--Bangert decomposition, and the strict maximum principle. These are reported ingredients, not supplied proofs. With the convention that the empty set has dimension $-\infty$, the stronger bound would require this aggregate singular set to be empty for $N\le7$.

For the second question, in the standard compact Riemannian interpretation let $M$ be closed and $0\le q\le\dim M$. For real homology $h\in H_q(M;\R)$ put
\[
 \|h\|_{\mathrm{st}}=
 \inf\{\mathbf M(T):\partial T=0,\ [T]=h\},\qquad
 B_{\mathrm{st}}=\{h:\|h\|_{\mathrm{st}}\le1\},
\]
where the infimum is over the appropriate finite-mass real currents and $\mathbf M$ is Riemannian mass, dual to the comass bound on test forms. The current class and compactness assumptions are not specified in the extract. In the usual stable-norm setting the aid records positivity, absolute homogeneity and the triangle inequality, so $B_{\mathrm{st}}$ is convex, centrally symmetric, and compact in the finite-dimensional homology space. The broad question concerns its finer geometry and regularity; elementary convexity alone does not settle it. At a nonzero class, differentiability of the norm can equivalently be expressed as uniqueness of its supporting dual functional. The aid also names a known Auer--Bangert codimension-one differentiability theorem, but gives no exact hypotheses or conclusion from which differentiability at every class could be asserted. The general-$q$ program is left open in its assessment.
\subsection{Short English statement}
Two concatenated questions concern singular sets of minimizing currents and the geometry of stable-norm balls. The aid reports a stronger aggregate singular-set bound and elementary stable-ball properties, while its precise singular-set definition and cited differentiability theorem remain insufficiently specified for independent verification.
\subsection{Sources}
\begin{itemize}
\item Local UnsolvedMath registry, permanent ID 20002485, source code catalogue entry; created and updated 14 August 2026. Canonical \texttt{aim-pdes-notes.json}, zero-based \texttt{notes[96]}, tag \texttt{problem}, background/extracted item 8 followed by 29.
\item AIM ``New connections between dynamical systems and PDE's'': \url{https://aimath.org/WWN/dynpde/dynpde.pdf}. The supplied aid identifies the concatenated questions as original Problems 28 and 29, the first attributed to Franz Auer. The verification flags concatenation, and the registry marks the wording exact. Original remarks and literature field are empty.
\item Research attempt 1: \texttt{partial\_result}, problem status at run open, rigor label \texttt{proved}, novelty confidence low. It names Chodosh--Mantoulidis--Schulze and Auer--Bangert, without full paper titles, dates, theorem/lemma statements, or authoritative bibliographic links.
\item Ulam AI / UnsolvedMath Contributors, CC BY 4.0. Machine-generated research aid, not peer reviewed; independently verify before citation.
\end{itemize}
\end{document}
